% ---------- 摘要（约300字） ----------
\begin{chineseabstract}

本课题为《微机原理》课程大作业题目3：设计并实现一套基于 Wi-Fi 信道状态信息（CSI）扰动特征的多监测点监控系统。系统由 3 个发送节点（2 块经典 ESP32、1 块 ESP32-C3）与 1 个 ESP32-S3 接收节点组成：发送节点以约 100~Hz 持续发包，接收节点用 esp-radar 组件按源 MAC 将 CSI 分发到各监测点链路，逐链路完成滑动窗均值/方差/峰峰值统计与滞回状态机判决，经 2~Mbps UART 二进制协议（帧同步 AA~55、CRC-16/CCITT-FALSE 校验、5 类帧）上报 PC 上位机（PySide6+pyqtgraph），实现区域状态卡片、实时特征曲线、事件日志与 CSV 记录/回放；标定阈值与链路表经 NVS 断电保存。台架实测：协议自测 20/20 通过；单链路 60~s 连续采集 546 帧无断流、包间隔中位数 110~ms；60~s 心跳 59 帧、特征帧 533 帧、零 CRC 错；双链路标定阈值 681/20367（centi-var）并经 TX 复位静默窗验证事件 MAC 归因无误；断电重启后链路表与阈值 1.26~s 内从 NVS 恢复；上位机冒烟测试 17/17 通过。系统功能链路已在台架全部打通，空间部署后的现场指标验收待执行。

\chinesekeyword{WiFi CSI；ESP32；多节点；扰动检测；UART}

\end{chineseabstract}
